\section{Experimental Setup}
\label{sec:experiments}

\subsection{Controlled SFT Injection}
\label{sec:injection}

We construct a controlled contamination testbed spanning five independent model families from five labs, trained on independent corpora with distinct pre-training recipes: Qwen2.5 \citep{qwen2024qwen25} (0.5B, 1.5B, 3B, 7B, 14B), Phi-4 \citep{abdin2024phi4} (14.7B), OLMo-2 \citep{olmo2team2025} (7B), Gemma-2 \citep{gemmateam2024gemma2} (9B), and Falcon3 \citep{tiiuae2024falcon3} (7B). All models are fine-tuned with a unified LoRA configuration \citep{hu2022lora} ($r=16$, $\alpha=32$, dropout$=0$, targeting all attention and multilayer perceptron (MLP) projections; lr$=5\times10^{-5}$, bf16, 1 epoch, batch size$=1$, gradient accumulation$=16$, warmup ratio$=0.03$; trainable parameters range from 0.45\% of Phi-4 14.7B to 2.10\% of Qwen2.5-0.5B; convergence in Appendix~\ref{sec:training_convergence}) on a general-purpose SFT corpus into which benchmark items are injected at five dosage levels: d000 (0\%, clean control), d005 (5\%), d025 (25\%), d050 (50\%), and d100 (100\%). Contaminated items augment (not replace) the base SFT corpus, yielding 50.0k (d000), 50.8k (d005), 54.1k (d025), 58.3k (d050), and 66.5k (d100) total training examples ($+$33\% volume at d100 vs.\ d000).

Qwen2.5 serves as the primary within-family series (five scale points, 0.5B--14B). Models at $\leq$3B are trained with two seeds at four dosage levels; 7B and 14B at all five dosages. The four cross-family models are each trained at five dosages with a single seed, serving as consistency checks rather than independent confirmatory analyses. This yields over 55 checkpoints spanning nine sizes across five families. We exclude 0.5B from MC IRT due to position bias (A-option rate: 46\% at d000 to 98\% at d100; \S\ref{sec:signal_failure}; details in Appendix~\ref{sec:exclusion_details}). We additionally validate our dose-response profiles on a held-out sixth family, LLaMA-3.1-8B \citep{grattafiori2024llama3} (Meta), which was not used during profile construction. LLaMA-3.1 is trained at three dosage levels (d000, d025, d100) to verify that Algorithm~\ref{alg:profile-classification} generalizes to unseen model families.

\subsection{Benchmarks, Domains, and Evaluation}
\label{sec:benchmarks}

We evaluate on three benchmarks spanning distinct cognitive demands and response formats.

\textbf{MMLU} \citep{hendrycks2021mmlu} (5-shot): 14,042 multiple-choice items across 57 subjects, grouped into 8 cognitive domains (Math, Science, CS \& Engineering, Medical, Humanities, Social Sciences, Business \& Law, Other). This stratification reveals within-benchmark heterogeneity that aggregate scoring obscures (\S\ref{sec:results}).

\textbf{ARC-Challenge} \citep{clark2018arc} (25-shot): 1,172 multiple-choice science reasoning items.

\textbf{GSM8K} \citep{cobbe2021gsm8k}: 1,319 open-ended math problems whose free-form generation format avoids multiple-choice artifacts. Gemma-2 is excluded from GSM8K cross-benchmark analysis because MC-format SFT destroys its chain-of-thought (CoT) reasoning: the d000 baseline drops to 17.8\% (vs.\ 73--84\% for other families), consistent with CoT collapse rather than a contamination effect.

All training runs were performed on NVIDIA RTX 4090D (24\,GB) GPUs; evaluation used NVIDIA RTX PRO 6000 (96\,GB) GPUs. Each LoRA fine-tuning run required 2--6 hours on a single GPU depending on model size and dosage; total compute across all 55+ checkpoints was approximately 200 GPU-hours.

All multiple-choice evaluations use logprob-based scoring under an aligned instruction-template format with fixed few-shot counts; all 75 model--dosage configurations are evaluated under this identical protocol.

We assess calibration quality via mean absolute error (MAE) (calibrated vs.\ ground-truth clean accuracy), Kendall $\tau$ (rank preservation), credible interval coverage, posterior predictive $p$-value (PPP), $\gamma$ direction ($P(\gamma\!<\!0 \mid \text{data})$), and split-$\hat{R}$ convergence.

\subsection{Baselines}
\label{sec:baselines}

We compare against seven baselines (Table~\ref{tab:baselines}): \textbf{B0-Raw} (no adjustment), \textbf{B1-Remove\&Rescore} (drop flagged items; B1b uses RRF-filtered items), \textbf{B2-Single2PL+$\gamma$} (aggregate $\gamma$ without domain stratification), \textbf{B3-Vanilla2PL} ($\gamma=0$, isolating IRT ability estimation), \textbf{B4-LinearRegression} (residuals on exposure scores), \textbf{B5-Emperor} (mitigation strategies from \citet{sun2025emperor}; B5b uses the best-performing variant), and \textbf{B6-DCR} (our reimplementation following \citealp{xu2025dcr}).

\begin{table}[t]
\centering
\caption{MMLU method comparison on a single random subsample of 300 items (Qwen2.5, $J$=16). ARC and GSM8K baselines in Appendix~\ref{sec:full_baselines}; ARC IRT diagnostics in Appendix~\ref{sec:arc_irt}.}
\label{tab:baselines}
\footnotesize
\begin{tabular}{lcc}
\toprule
Method & MAE & $\tau$ \\
\midrule
B0: Raw Accuracy & \textbf{0.006} & 0.767 \\
B1b: Remove \& Rescore & 0.026 & \underline{0.783} \\
B3: Vanilla 2PL & 0.008 & 0.731 \\
B4: Linear Regression & 0.010 & 0.751 \\
B5b: Emperor & \textbf{0.006} & \textbf{0.789} \\
B6: DCR (reimpl.) & 0.108 & 0.654 \\
ContamCal (2PL+$\gamma$) & 0.009 & 0.731 \\
\bottomrule
\end{tabular}
\end{table}